\section{AI 产品组织：Taste 与两段式}

当产品从流程设计转为上下文设计，组织方式也必须变。张月光说，决策边界变模糊，Taste 变得重要，工作从线性变成两段式。这里的 Taste，不是主观审美，而是判断模型输出好坏、用户体验边界和产品方向是否值得押注的能力。

\begin{figure}[H]
\centering
\includegraphics[width=0.96\textwidth]{ai-product-org-two-stage.png}
\caption{AI 产品组织：线性工作变两段式。}
\end{figure}

\begin{knowledgebox}{读图：先验证能力，再封装产品}
图中第一段是模型可行性：context、prompt、工具、模型组合、评测和 taste 判断；第二段才是产品封装：流程、UI、商业模式和用户心智。传统线性工作流把需求定义放在最前面，AI Native 团队则要先确认模型能力真实存在。
\end{knowledgebox}

\subsection{Taste 为什么重要}

AI 输出有不确定性，用户输入也不确定。产品经理不能只说“这里放一个按钮”，而要判断输出是否足够好、是否可控、是否有用户愿意持续使用、是否能成为单向门。这个判断往往无法完全被指标替代，因为早期产品还没有稳定指标。

\begin{knowledgebox}{术语消化：AI 产品组织}
\begin{tabularx}{\textwidth}{p{0.24\textwidth}p{0.32\textwidth}X}
\toprule
术语 & 一句话解释 & 在本期中的作用 \\
\midrule
Taste & 对输出质量、产品方向和用户体验的判断力 & 决策边界模糊时的核心能力。 \\
Context Design & 设计模型可见的信息、工具和约束 & AI Native 产品的中间层。 \\
Prompt Practice & 构造输入和任务描述的方法 & 只是 context design 的一部分。 \\
Two-stage Work & 先验证模型能力，再封装产品流程 & 替代传统线性协作。 \\
Evaluation & 判断输出是否稳定、可用和可迭代 & 连接 taste 与工程。 \\
\bottomrule
\end{tabularx}
\end{knowledgebox}

\subsection{本章小结}

AI 产品组织的关键，是让产品、设计、工程和模型能力在同一个反馈环里工作。没有 taste，团队容易被 demo 迷惑；没有工程验证，taste 又会变成玄学。

